% -*- coding: utf-8 -*-
% !TEX program = xelatex
\documentclass[plain,noanswer,medmath]{../../template/jnuexam}
\usepackage{../../template/kysx}

\begin{document}


\examtitle{2012}{1}


\makepart{选择题}{1～8小题，每小题4分，共32分}


\begin{problem}
曲线$y = \frac{x^2 + x}{x^2 - 1}$渐近线的条数为\pickout{}
\begin{abcd}
\item $0$．
\item $1$．
\item $2$．
\item $3$．
\end{abcd}
\end{problem}


\begin{problem}
设函数$f(x) = (\e^x - 1)(\e^{2x} - 2) \cdots(\e^{nx} - n)$，其中$n$为正整数，则$f'(0) =$\pickout{}
\begin{abcd}
\item $(- 1)^{n - 1}(n - 1)!$．
\item $(- 1)^n(n - 1)!$．
\item $(- 1)^{n - 1} n!$．
\item $(- 1)^n n!$．
\end{abcd}
\end{problem}


\begin{problem}
如果$f(x,y)$在$\left(0,0 \right)$处连续，那么下列命题正确的是\pickout{}
\begin{abcd}
\item 若极限$\lim_{\begin{subarray}{l}
  x \to 0 \\
  y \to 0
\end{subarray}} \frac{f(x,y)}{|x| + |y|}$存在，则$f(x,y)$在$(0,0)$处可微．
\item 若极限$\lim_{\begin{subarray}{l}
  x \to 0 \\
  y \to 0
\end{subarray}} \frac{f(x,y)}{x^2 + y^2}$存在，则$f(x,y)$在$(0,0)$处可微．
\item 若$f(x,y)$在$(0,0)$处可微，则极限$\lim_{\begin{subarray}{l}
  x \to 0 \\
  y \to 0
\end{subarray}} \frac{f(x,y)}{|x| + |y|}$存在．
\item 若$f(x,y)$在$(0,0)$处可微，则极限$\lim_{\begin{subarray}{l}
  x \to 0 \\
  y \to 0
\end{subarray}} \frac{f(x,y)}{x^2 + y^2}$存在．
\end{abcd}
\end{problem}


\begin{problem}
设$I_k = \int_0^{k\pi} \e^{x^{2}} \sin x\dx$（$k = 1,2,3$），则有\pickout{}
\begin{abcd}
\item $I_1 < I_2 < I_3$．
\item $I_3 < I_2 < I_1$．
\item $I_2 < I_3 < I_1$．
\item $I_2 < I_1 < I_3$．
\end{abcd}
\end{problem}


\begin{problem}
设$\alpha_1 = \begin{pmatrix}
  0 \\
  0 \\
  c_1
\end{pmatrix},\alpha_2 = \begin{pmatrix}
  0 \\
  1 \\
  c_2
\end{pmatrix},\alpha_3 = \begin{pmatrix}
  1 \\
  -1 \\
  c_3
\end{pmatrix},\alpha_4 = \begin{pmatrix}
  -1 \\
  1 \\
  c_4
\end{pmatrix}$其中$c_1,c_2,c_3,c_4$为任意常数，则下列向量组线性相关的是\pickout{}
\begin{abcd}
\item $\alpha_1,\alpha_2,\alpha_3$．
\item $\alpha_1,\alpha_2,\alpha_4$．
\item $\alpha_1,\alpha_3,\alpha_4$．
\item $\alpha_2,\alpha_3,\alpha_4$．
\end{abcd}
\end{problem}


\begin{problem}
设$A$为$3$阶矩阵，$P$为$3$阶可逆矩阵，且$P^{-1}AP = \begin{pmatrix}
  1 &   & \\
    & 1 & \\
    &   & 2
\end{pmatrix}$，$P = \left(\alpha_1,\alpha_2,\alpha_3 \right)$，
$Q = \left(\alpha_1 + \alpha_2,\alpha_2,\alpha_3 \right)$则$Q^{-1}AQ=$\pickout{}
\begin{abcd}
\item $\begin{pmatrix}
  1 &   & \\
    & 2 & \\
    &   & 1
\end{pmatrix}$．
\item $\begin{pmatrix}
  1 &   & \\
    & 1 & \\
    &   & 2
\end{pmatrix}$．
\item $\begin{pmatrix}
  2 &   & \\
    & 1 & \\
    &   & 2
\end{pmatrix}$．
\item $\begin{pmatrix}
  2 &   & \\
    & 2 & \\
    &   & 1
\end{pmatrix}$．
\end{abcd}
\end{problem}


\begin{problem}
设随机变量$x$与$y$相互独立，且分别服从参数为$1$与参数为$4$的指数分布，则$P\{x < y\}=$\pickout{}
\begin{abcd}
\item $\frac{1}{5}$．
\item $\frac{1}{3}$．
\item $\frac{2}{3}$．
\item $\frac{4}{5}$．
\end{abcd}
\end{problem}


\begin{problem}
将长度为$1m$的木棒随机地截成两段，则两段长度的相关系数为\pickout{}
\begin{abcd}
\item $1$．
\item $\frac{1}{2}$．
\item $- \frac{1}{2}$．
\item $- 1$．
\end{abcd}
\end{problem}


\makepart{填空题}{9～14小题，每小题4分，共24分}


\begin{problem}
若函数$f(x)$满足方程$f''(x) + f'(x) - 2f(x) = 0$及$f'(x) + f(x) = 2\e^x$，则$f(x) =$ \fillin{}．
\end{problem}


\begin{problem}
$\int_0^2 x\sqrt{2x - x^2} \dx =$ \fillin{}．
\end{problem}


\begin{problem}
$\left.\grad\left(xy + \frac{z}{y} \right) \right|_{(2,1,1)} =$ \fillin{}．
\end{problem}


\begin{problem}
设$\Sigma = \left\{\left(x,y,z \right)\left| x + y + z = 1,x \ge 0,y \ge 0,z \ge 0 \right. \right\}$，
则$\iint_{\Sigma}y^2 \dS =$ \fillin{}．
\end{problem}


\begin{problem}
设$\alpha$为$3$维单位列向量，$E$为$3$阶单位矩阵，则矩阵$E - \alpha \alpha^T$的秩为 \fillin{}．
\end{problem}


\begin{problem}
设$A,B,C$是随机事件，$A$与$C$互不相容，$P(AB) = \frac{1}{2}$, $P(C) = \frac{1}{3}$，
则$P(AB|\overline{C}) =$ \fillin{}．
\end{problem}


\makepart{解答题}{15～23小题，共94分}


\begin{problem}[points=10]
证明：$x\ln \frac{1 + x}{1 - x} + \cos x \ge 1 + \frac{x^2}{2}$（$- 1 < x < 1$）．
\end{problem}


\begin{problem}[points=10]
求$f\left(x,y \right) = x\e^{- \frac{x^2 + y^2}{2}}$的极值．
\end{problem}


\begin{problem}[points=10]
求幂级数$\sum_{n = 0}^{\infty}\frac{4n^2 + 4n + 3}{2n + 1}x^{2n}$的收敛域及和函数．
\end{problem}


\begin{problem}[points=10]
已知曲线$L:\begin{cases}
  x = f(t) \\
  y = \cos t
\end{cases}$($0 \le t < \frac{\pi}{2}$)，其中函数$f(t)$具有连续导数，
且$f(0) = 0$，$f(t) > 0\left(0 < t < \frac{\pi}{2} \right)$．
若曲线$L$的切线与$x$轴的交点到切点的距离恒为$1$，求函数$f(t)$的表达式，
并求以曲线$L$及$x$轴与$y$轴为边界的区域的面积．
\end{problem}


\begin{problem}[points=10]
已知$L$是第一象限中从点$\left(0,0 \right)$沿圆周$x^2 + y^2 = 2x$到点$\left(2,0 \right)$，
再沿圆周\goodbreak$x^2 + y^2 = 4$到点$\left(0,2 \right)$的曲线段，计算曲线积分
$$I=\int_L 3x^2 y\dx + \left(x^3 + x - 2y \right)\dy$$
\end{problem}


\begin{problem}[points=11]
设$A = \begin{pmatrix}
  1 & a & 0 & 0 \\
  0 & 1 & a & 0 \\
  0 & 0 & 1 & a \\
  a & 0 & 0 & 1
\end{pmatrix}$，$b = \begin{pmatrix}
  1 \\
  -1 \\
  0 \\
  0
\end{pmatrix}$．
\begin{enumerate}
\item 计算行列式$|A|$．
\item 当实数$a$为何值时，方程组$Ax = b$有无穷多解，并求其通解．
\end{enumerate}
\end{problem}


\begin{problem}[points=11]
设$A = \begin{pmatrix}
   1 & 0 & 1 \\
   0 & 1 & 1 \\
  -1 & 0 & a \\
   0 & a & -1
\end{pmatrix}$，二次型$f(x_1,x_2,x_3) = x^T A^T Ax$的秩为$2$．\par
\begin{enumerate*}
\item 求实数$a$的值；\item 求正交变换$x = Qy$将二次型$f$化为标准型．
\end{enumerate*}
\end{problem}


\begin{problem}[points=11]
设二维离散型随机变量$(X,Y)$的概率分布为
\[\begin{array}{|c|c|c|c|}
\hline
  \diagboxtwo{X}{Y} & 0    & 1   & 2    \\
\hline
       0            & 1/4  & 0   & 1/4  \\
\hline
       1            & 0    & 1/3 & 0    \\
\hline
       2            & 1/12 & 0   & 1/12 \\
\hline
\end{array}\]
\begin{enumerate*}
\item 求$P(X = 2Y)$；\item 求$\Cov \left(X - Y,Y \right)$．
\end{enumerate*}

\end{problem}

\begin{problem}[points=11]
设随机变量$X$与$Y$相互独立且分别服从正态分布$N\left(\mu ,\sigma^2 \right)$与
$N\left(\mu ,2\sigma^2 \right)$，其中$\sigma$是未知参数且$\sigma > 0$，设$Z = X - Y$,
\begin{enumerate}
\item 求$Z$的概率密度$f\left(z, \sigma^2 \right)$；
\item 设$Z_1,Z_2, \cdots Z_n$为来自总体$Z$的简单随机样本，求$\sigma^2$的最大似然估计量$\hat{\sigma}^2$；
\item 证明$\hat{\sigma}^2$为$\sigma^2$的无偏估计量．
\end{enumerate}
\end{problem}


\end{document}
